\clearpage\phantomsection\label{couverture-programme}\pdfbookmark[0]{Repérage du programme}{bm-couverture-programme}\gdef\otChapNum{0}\pagestyle{otFront}
\otPageTitre{Repérage du programme}{CM1 · Programme 2025 du cycle 3}
Cette matrice relie les objectifs d’apprentissage du programme aux
exercices du livre. Les cours complètent ces repères, notamment pour
le vocabulaire, les représentations et les méthodes de construction.\par\bigskip
Les six domaines du programme sont couverts : nombres, calcul et
résolution de problèmes ; grandeurs et mesures ; espace et géométrie ;
organisation et gestion de données et probabilités ; proportionnalité ;
initiation à la pensée informatique, travaillée dans les chapitres 5,
15, 16 et 18 ainsi que dans plusieurs défis.\par\bigskip
\begin{attention}Les niveaux de difficulté et les défis sont des choix
d’entraînement. Ils ne constituent pas une liste d’obligations
supplémentaires au programme officiel.\end{attention}\clearpage
\otPageTitre{Repères et exercices}{Chapitres 1 à 3}\fontsize{10}{13}\selectfont
\otIntertitre{1. Les nombres jusqu’au million}\begin{tabularx}{\linewidth}{@{}Yl@{}}\toprule
Classes, rangs et décompositions & \hyperref[exo-1-1]{1.1}, \hyperref[exo-1-2]{1.2}, \hyperref[exo-1-3]{1.3}, \hyperref[exo-1-6]{1.6}, \hyperref[exo-1-7]{1.7}\\[2mm]
Chiffre et nombre d’unités & \hyperref[exo-1-4]{1.4}, \hyperref[exo-1-12]{1.12}\\[2mm]
Comparer, encadrer, intercaler & \hyperref[exo-1-5]{1.5}, \hyperref[exo-1-9]{1.9}, \hyperref[exo-1-10]{1.10}, \hyperref[exo-1-13]{1.13}\\[2mm]
Demi-droite graduée et problèmes & \hyperref[exo-1-8]{1.8}, \hyperref[exo-1-11]{1.11}, \hyperref[exo-1-14]{1.14}, \hyperref[exo-1-15]{1.15}\\[2mm]
\bottomrule\end{tabularx}\par\vspace{4mm}
\otIntertitre{2. Additionner et soustraire}\begin{tabularx}{\linewidth}{@{}Yl@{}}\toprule
Techniques posées et contrôles & \hyperref[exo-2-1]{2.1}, \hyperref[exo-2-2]{2.2}, \hyperref[exo-2-7]{2.7}, \hyperref[exo-2-10]{2.10}\\[2mm]
Calcul mental et estimations & \hyperref[exo-2-3]{2.3}, \hyperref[exo-2-4]{2.4}, \hyperref[exo-2-5]{2.5}, \hyperref[exo-2-6]{2.6}, \hyperref[exo-2-15]{2.15}\\[2mm]
Problèmes parties-tout & \hyperref[exo-2-9]{2.9}, \hyperref[exo-2-11]{2.11}, \hyperref[exo-2-13]{2.13}\\[2mm]
Comparaisons et étapes multiples & \hyperref[exo-2-8]{2.8}, \hyperref[exo-2-12]{2.12}, \hyperref[exo-2-14]{2.14}\\[2mm]
\bottomrule\end{tabularx}\par\vspace{4mm}
\otIntertitre{3. Multiplier des entiers}\begin{tabularx}{\linewidth}{@{}Yl@{}}\toprule
Sens du produit et tables & \hyperref[exo-3-1]{3.1}, \hyperref[exo-3-2]{3.2}, \hyperref[exo-3-9]{3.9}\\[2mm]
Technique posée et estimation & \hyperref[exo-3-4]{3.4}, \hyperref[exo-3-6]{3.6}, \hyperref[exo-3-15]{3.15}\\[2mm]
Multiplier par 10, 100 et des dizaines & \hyperref[exo-3-3]{3.3}, \hyperref[exo-3-5]{3.5}, \hyperref[exo-3-7]{3.7}, \hyperref[exo-3-8]{3.8}\\[2mm]
Problèmes multiplicatifs et dénombrement & \hyperref[exo-3-10]{3.10}, \hyperref[exo-3-11]{3.11}, \hyperref[exo-3-12]{3.12}, \hyperref[exo-3-13]{3.13}, \hyperref[exo-3-14]{3.14}\\[2mm]
\bottomrule\end{tabularx}\par\vspace{4mm}
\clearpage
\otPageTitre{Repères et exercices}{Chapitres 4 à 6}\fontsize{10}{13}\selectfont
\otIntertitre{4. Diviser : partages et restes}\begin{tabularx}{\linewidth}{@{}Yl@{}}\toprule
Multiples, pairs et critères & \hyperref[exo-4-1]{4.1}, \hyperref[exo-4-2]{4.2}, \hyperref[exo-4-3]{4.3}\\[2mm]
Division posée et contrôle & \hyperref[exo-4-4]{4.4}, \hyperref[exo-4-6]{4.6}, \hyperref[exo-4-7]{4.7}, \hyperref[exo-4-10]{4.10}\\[2mm]
Quotient, reste et encadrements & \hyperref[exo-4-5]{4.5}, \hyperref[exo-4-8]{4.8}, \hyperref[exo-4-9]{4.9}\\[2mm]
Partages et problèmes de groupes & \hyperref[exo-4-11]{4.11}, \hyperref[exo-4-12]{4.12}, \hyperref[exo-4-13]{4.13}, \hyperref[exo-4-14]{4.14}, \hyperref[exo-4-15]{4.15}\\[2mm]
\bottomrule\end{tabularx}\par\vspace{4mm}
\otIntertitre{5. Égalités, inconnues et régularités}\begin{tabularx}{\linewidth}{@{}Yl@{}}\toprule
Sens du signe égal & \hyperref[exo-5-1]{5.1}\\[2mm]
Égalités à trous et nombres cachés & \hyperref[exo-5-2]{5.2}, \hyperref[exo-5-3]{5.3}, \hyperref[exo-5-7]{5.7}, \hyperref[exo-5-9]{5.9}, \hyperref[exo-5-10]{5.10}, \hyperref[exo-5-12]{5.12}\\[2mm]
Suites de nombres et de motifs & \hyperref[exo-5-5]{5.5}, \hyperref[exo-5-6]{5.6}, \hyperref[exo-5-8]{5.8}, \hyperref[exo-5-13]{5.13}, \hyperref[exo-5-14]{5.14}\\[2mm]
Programmes de calcul et lettres & \hyperref[exo-5-4]{5.4}, \hyperref[exo-5-11]{5.11}, \hyperref[exo-5-15]{5.15}\\[2mm]
\bottomrule\end{tabularx}\par\vspace{4mm}
\otIntertitre{6. Découvrir les fractions}\begin{tabularx}{\linewidth}{@{}Yl@{}}\toprule
Parts égales et vocabulaire & \hyperref[exo-6-1]{6.1}, \hyperref[exo-6-2]{6.2}, \hyperref[exo-6-3]{6.3}, \hyperref[exo-6-4]{6.4}\\[2mm]
Fractions supérieures à 1 & \hyperref[exo-6-5]{6.5}, \hyperref[exo-6-6]{6.6}, \hyperref[exo-6-13]{6.13}\\[2mm]
Encadrements et entiers cachés & \hyperref[exo-6-7]{6.7}, \hyperref[exo-6-8]{6.8}, \hyperref[exo-6-10]{6.10}\\[2mm]
Partages et fractions d’un tout & \hyperref[exo-6-9]{6.9}, \hyperref[exo-6-11]{6.11}, \hyperref[exo-6-12]{6.12}, \hyperref[exo-6-14]{6.14}, \hyperref[exo-6-15]{6.15}\\[2mm]
\bottomrule\end{tabularx}\par\vspace{4mm}
\clearpage
\otPageTitre{Repères et exercices}{Chapitres 7 à 9}\fontsize{10}{13}\selectfont
\otIntertitre{7. Utiliser et comparer les fractions}\begin{tabularx}{\linewidth}{@{}Yl@{}}\toprule
Demi-droite et bande unité & \hyperref[exo-7-1]{7.1}, \hyperref[exo-7-2]{7.2}, \hyperref[exo-7-9]{7.9}, \hyperref[exo-7-11]{7.11}\\[2mm]
Comparer des fractions & \hyperref[exo-7-3]{7.3}, \hyperref[exo-7-4]{7.4}, \hyperref[exo-7-10]{7.10}, \hyperref[exo-7-14]{7.14}\\[2mm]
Sommes et différences de parts & \hyperref[exo-7-6]{7.6}, \hyperref[exo-7-7]{7.7}, \hyperref[exo-7-8]{7.8}, \hyperref[exo-7-13]{7.13}\\[2mm]
Fraction d’une quantité et problèmes & \hyperref[exo-7-5]{7.5}, \hyperref[exo-7-12]{7.12}, \hyperref[exo-7-15]{7.15}\\[2mm]
\bottomrule\end{tabularx}\par\vspace{4mm}
\otIntertitre{8. Fractions décimales et virgule}\begin{tabularx}{\linewidth}{@{}Yl@{}}\toprule
Dixièmes, centièmes et unités & \hyperref[exo-8-1]{8.1}, \hyperref[exo-8-2]{8.2}, \hyperref[exo-8-3]{8.3}\\[2mm]
Décompositions décimales & \hyperref[exo-8-4]{8.4}, \hyperref[exo-8-5]{8.5}, \hyperref[exo-8-6]{8.6}, \hyperref[exo-8-13]{8.13}\\[2mm]
Écriture à virgule et rangs & \hyperref[exo-8-7]{8.7}, \hyperref[exo-8-8]{8.8}, \hyperref[exo-8-10]{8.10}, \hyperref[exo-8-12]{8.12}\\[2mm]
Monnaie, mesures et bons zéros & \hyperref[exo-8-9]{8.9}, \hyperref[exo-8-11]{8.11}, \hyperref[exo-8-14]{8.14}, \hyperref[exo-8-15]{8.15}\\[2mm]
\bottomrule\end{tabularx}\par\vspace{4mm}
\otIntertitre{9. Comparer et arrondir les décimaux}\begin{tabularx}{\linewidth}{@{}Yl@{}}\toprule
Comparer rang par rang & \hyperref[exo-9-1]{9.1}, \hyperref[exo-9-2]{9.2}, \hyperref[exo-9-3]{9.3}, \hyperref[exo-9-10]{9.10}\\[2mm]
Encadrer et intercaler & \hyperref[exo-9-4]{9.4}, \hyperref[exo-9-7]{9.7}, \hyperref[exo-9-13]{9.13}\\[2mm]
Demi-droite et graduations fines & \hyperref[exo-9-6]{9.6}, \hyperref[exo-9-9]{9.9}\\[2mm]
Partie entière, arrondis et problèmes & \hyperref[exo-9-5]{9.5}, \hyperref[exo-9-8]{9.8}, \hyperref[exo-9-11]{9.11}, \hyperref[exo-9-12]{9.12}, \hyperref[exo-9-14]{9.14}, \hyperref[exo-9-15]{9.15}\\[2mm]
\bottomrule\end{tabularx}\par\vspace{4mm}
\clearpage
\otPageTitre{Repères et exercices}{Chapitres 10 à 12}\fontsize{10}{13}\selectfont
\otIntertitre{10. Calculer avec les décimaux}\begin{tabularx}{\linewidth}{@{}Yl@{}}\toprule
Poser avec les virgules alignées & \hyperref[exo-10-1]{10.1}, \hyperref[exo-10-2]{10.2}, \hyperref[exo-10-7]{10.7}, \hyperref[exo-10-9]{10.9}\\[2mm]
Calcul mental et compléments & \hyperref[exo-10-3]{10.3}, \hyperref[exo-10-6]{10.6}, \hyperref[exo-10-8]{10.8}, \hyperref[exo-10-10]{10.10}\\[2mm]
Multiplier par 10 et par un entier & \hyperref[exo-10-4]{10.4}, \hyperref[exo-10-5]{10.5}\\[2mm]
Monnaie, mesures et contrôles & \hyperref[exo-10-11]{10.11}, \hyperref[exo-10-12]{10.12}, \hyperref[exo-10-13]{10.13}, \hyperref[exo-10-14]{10.14}, \hyperref[exo-10-15]{10.15}\\[2mm]
\bottomrule\end{tabularx}\par\vspace{4mm}
\otIntertitre{11. Longueurs et périmètres}\begin{tabularx}{\linewidth}{@{}Yl@{}}\toprule
Unités et conversions raisonnées & \hyperref[exo-11-1]{11.1}, \hyperref[exo-11-2]{11.2}, \hyperref[exo-11-3]{11.3}, \hyperref[exo-11-9]{11.9}\\[2mm]
Estimations et références & \hyperref[exo-11-4]{11.4}\\[2mm]
Périmètres de polygones & \hyperref[exo-11-5]{11.5}, \hyperref[exo-11-6]{11.6}, \hyperref[exo-11-7]{11.7}, \hyperref[exo-11-8]{11.8}\\[2mm]
Figures composées et problèmes & \hyperref[exo-11-10]{11.10}, \hyperref[exo-11-11]{11.11}, \hyperref[exo-11-12]{11.12}, \hyperref[exo-11-13]{11.13}, \hyperref[exo-11-14]{11.14}, \hyperref[exo-11-15]{11.15}\\[2mm]
\bottomrule\end{tabularx}\par\vspace{4mm}
\otIntertitre{12. Masses et contenances}\begin{tabularx}{\linewidth}{@{}Yl@{}}\toprule
Unités de masse et références & \hyperref[exo-12-2]{12.2}, \hyperref[exo-12-4]{12.4}, \hyperref[exo-12-6]{12.6}, \hyperref[exo-12-7]{12.7}\\[2mm]
Unités de contenance et références & \hyperref[exo-12-3]{12.3}, \hyperref[exo-12-8]{12.8}, \hyperref[exo-12-13]{12.13}\\[2mm]
Conversions raisonnées & \hyperref[exo-12-1]{12.1}, \hyperref[exo-12-5]{12.5}, \hyperref[exo-12-9]{12.9}, \hyperref[exo-12-10]{12.10}\\[2mm]
Problèmes de masses et de volumes utiles & \hyperref[exo-12-11]{12.11}, \hyperref[exo-12-12]{12.12}, \hyperref[exo-12-14]{12.14}, \hyperref[exo-12-15]{12.15}\\[2mm]
\bottomrule\end{tabularx}\par\vspace{4mm}
\clearpage
\otPageTitre{Repères et exercices}{Chapitres 13 à 15}\fontsize{10}{13}\selectfont
\otIntertitre{13. Aires et angles : comparer}\begin{tabularx}{\linewidth}{@{}Yl@{}}\toprule
Comparer des surfaces & \hyperref[exo-13-1]{13.1}, \hyperref[exo-13-3]{13.3}, \hyperref[exo-13-9]{13.9}, \hyperref[exo-13-15]{13.15}\\[2mm]
Pavages et centimètre carré & \hyperref[exo-13-2]{13.2}, \hyperref[exo-13-4]{13.4}, \hyperref[exo-13-8]{13.8}, \hyperref[exo-13-11]{13.11}, \hyperref[exo-13-12]{13.12}, \hyperref[exo-13-13]{13.13}\\[2mm]
Angles, sommet et notation & \hyperref[exo-13-5]{13.5}, \hyperref[exo-13-6]{13.6}, \hyperref[exo-13-10]{13.10}\\[2mm]
Angle droit et classements & \hyperref[exo-13-7]{13.7}, \hyperref[exo-13-14]{13.14}\\[2mm]
\bottomrule\end{tabularx}\par\vspace{4mm}
\otIntertitre{14. Horaires et durées}\begin{tabularx}{\linewidth}{@{}Yl@{}}\toprule
Lecture de l’heure & \hyperref[exo-14-1]{14.1}, \hyperref[exo-14-2]{14.2}, \hyperref[exo-14-3]{14.3}\\[2mm]
Conversions heures-minutes & \hyperref[exo-14-4]{14.4}, \hyperref[exo-14-9]{14.9}, \hyperref[exo-14-10]{14.10}\\[2mm]
Durées entre deux instants & \hyperref[exo-14-5]{14.5}, \hyperref[exo-14-6]{14.6}, \hyperref[exo-14-13]{14.13}\\[2mm]
Horaires de début et de fin & \hyperref[exo-14-7]{14.7}, \hyperref[exo-14-8]{14.8}, \hyperref[exo-14-11]{14.11}, \hyperref[exo-14-12]{14.12}, \hyperref[exo-14-14]{14.14}, \hyperref[exo-14-15]{14.15}\\[2mm]
\bottomrule\end{tabularx}\par\vspace{4mm}
\otIntertitre{15. Droites, cercles et constructions}\begin{tabularx}{\linewidth}{@{}Yl@{}}\toprule
Notations, codes et instruments & \hyperref[exo-15-1]{15.1}, \hyperref[exo-15-2]{15.2}, \hyperref[exo-15-9]{15.9}\\[2mm]
Perpendiculaires et parallèles & \hyperref[exo-15-3]{15.3}, \hyperref[exo-15-4]{15.4}, \hyperref[exo-15-6]{15.6}, \hyperref[exo-15-7]{15.7}, \hyperref[exo-15-14]{15.14}\\[2mm]
Cercle, disque et distances & \hyperref[exo-15-5]{15.5}, \hyperref[exo-15-8]{15.8}, \hyperref[exo-15-10]{15.10}, \hyperref[exo-15-12]{15.12}\\[2mm]
Programmes de construction & \hyperref[exo-15-11]{15.11}, \hyperref[exo-15-13]{15.13}, \hyperref[exo-15-15]{15.15}\\[2mm]
\bottomrule\end{tabularx}\par\vspace{4mm}
\clearpage
\otPageTitre{Repères et exercices}{Chapitres 16 à 18}\fontsize{10}{13}\selectfont
\otIntertitre{16. Triangles et quadrilatères}\begin{tabularx}{\linewidth}{@{}Yl@{}}\toprule
Triangles particuliers et codes & \hyperref[exo-16-1]{16.1}, \hyperref[exo-16-2]{16.2}, \hyperref[exo-16-5]{16.5}\\[2mm]
Constructions au compas & \hyperref[exo-16-3]{16.3}, \hyperref[exo-16-7]{16.7}, \hyperref[exo-16-10]{16.10}\\[2mm]
Quadrilatères usuels et propriétés & \hyperref[exo-16-4]{16.4}, \hyperref[exo-16-13]{16.13}, \hyperref[exo-16-15]{16.15}\\[2mm]
Rectangles, carrés et losanges construits & \hyperref[exo-16-6]{16.6}, \hyperref[exo-16-8]{16.8}, \hyperref[exo-16-9]{16.9}, \hyperref[exo-16-11]{16.11}, \hyperref[exo-16-12]{16.12}, \hyperref[exo-16-14]{16.14}\\[2mm]
\bottomrule\end{tabularx}\par\vspace{4mm}
\otIntertitre{17. La symétrie axiale}\begin{tabularx}{\linewidth}{@{}Yl@{}}\toprule
Axes et pliages & \hyperref[exo-17-1]{17.1}, \hyperref[exo-17-2]{17.2}, \hyperref[exo-17-3]{17.3}, \hyperref[exo-17-8]{17.8}\\[2mm]
Symétriques de points et segments & \hyperref[exo-17-4]{17.4}, \hyperref[exo-17-6]{17.6}, \hyperref[exo-17-9]{17.9}, \hyperref[exo-17-10]{17.10}\\[2mm]
Figures complétées sur quadrillage & \hyperref[exo-17-5]{17.5}, \hyperref[exo-17-7]{17.7}, \hyperref[exo-17-11]{17.11}, \hyperref[exo-17-12]{17.12}, \hyperref[exo-17-14]{17.14}\\[2mm]
Propriétés conservées & \hyperref[exo-17-13]{17.13}, \hyperref[exo-17-15]{17.15}\\[2mm]
\bottomrule\end{tabularx}\par\vspace{4mm}
\otIntertitre{18. Solides, patrons et espace}\begin{tabularx}{\linewidth}{@{}Yl@{}}\toprule
Solides usuels et vocabulaire & \hyperref[exo-18-1]{18.1}, \hyperref[exo-18-2]{18.2}, \hyperref[exo-18-3]{18.3}, \hyperref[exo-18-4]{18.4}, \hyperref[exo-18-5]{18.5}\\[2mm]
Faces, arêtes et descriptions & \hyperref[exo-18-10]{18.10}, \hyperref[exo-18-12]{18.12}\\[2mm]
Patrons du cube & \hyperref[exo-18-6]{18.6}, \hyperref[exo-18-7]{18.7}, \hyperref[exo-18-8]{18.8}, \hyperref[exo-18-11]{18.11}, \hyperref[exo-18-15]{18.15}\\[2mm]
Assemblages et déplacements codés & \hyperref[exo-18-9]{18.9}, \hyperref[exo-18-13]{18.13}, \hyperref[exo-18-14]{18.14}\\[2mm]
\bottomrule\end{tabularx}\par\vspace{4mm}
\clearpage
\otPageTitre{Repères et exercices}{Chapitres 19 à 21}\fontsize{10}{13}\selectfont
\otIntertitre{19. La proportionnalité}\begin{tabularx}{\linewidth}{@{}Yl@{}}\toprule
Reconnaître la proportionnalité & \hyperref[exo-19-1]{19.1}, \hyperref[exo-19-10]{19.10}, \hyperref[exo-19-15]{19.15}\\[2mm]
Fois plus et fois moins & \hyperref[exo-19-2]{19.2}, \hyperref[exo-19-3]{19.3}, \hyperref[exo-19-4]{19.4}, \hyperref[exo-19-5]{19.5}, \hyperref[exo-19-6]{19.6}, \hyperref[exo-19-7]{19.7}\\[2mm]
Situations régulières & \hyperref[exo-19-8]{19.8}, \hyperref[exo-19-9]{19.9}, \hyperref[exo-19-11]{19.11}, \hyperref[exo-19-12]{19.12}\\[2mm]
Pièges et comparaisons & \hyperref[exo-19-13]{19.13}, \hyperref[exo-19-14]{19.14}\\[2mm]
\bottomrule\end{tabularx}\par\vspace{4mm}
\otIntertitre{20. Organiser et lire des données}\begin{tabularx}{\linewidth}{@{}Yl@{}}\toprule
Tableaux simples et doubles & \hyperref[exo-20-1]{20.1}, \hyperref[exo-20-2]{20.2}, \hyperref[exo-20-3]{20.3}, \hyperref[exo-20-8]{20.8}, \hyperref[exo-20-11]{20.11}, \hyperref[exo-20-12]{20.12}\\[2mm]
Diagrammes en barres & \hyperref[exo-20-4]{20.4}, \hyperref[exo-20-6]{20.6}, \hyperref[exo-20-10]{20.10}\\[2mm]
Courbes et repères & \hyperref[exo-20-5]{20.5}, \hyperref[exo-20-7]{20.7}, \hyperref[exo-20-9]{20.9}, \hyperref[exo-20-13]{20.13}, \hyperref[exo-20-14]{20.14}\\[2mm]
Enquêtes et données réelles & \hyperref[exo-20-15]{20.15}\\[2mm]
\bottomrule\end{tabularx}\par\vspace{4mm}
\otIntertitre{21. Premiers pas vers le hasard}\begin{tabularx}{\linewidth}{@{}Yl@{}}\toprule
Expériences et issues & \hyperref[exo-21-1]{21.1}, \hyperref[exo-21-3]{21.3}\\[2mm]
Vocabulaire des chances & \hyperref[exo-21-2]{21.2}, \hyperref[exo-21-4]{21.4}, \hyperref[exo-21-5]{21.5}\\[2mm]
Équiprobabilité & \hyperref[exo-21-7]{21.7}, \hyperref[exo-21-8]{21.8}, \hyperref[exo-21-13]{21.13}, \hyperref[exo-21-14]{21.14}\\[2mm]
Comparer des évènements & \hyperref[exo-21-6]{21.6}, \hyperref[exo-21-9]{21.9}, \hyperref[exo-21-10]{21.10}, \hyperref[exo-21-11]{21.11}, \hyperref[exo-21-12]{21.12}, \hyperref[exo-21-15]{21.15}\\[2mm]
\bottomrule\end{tabularx}\par\vspace{4mm}
\clearpage
